\documentclass[11pt]{article}
\usepackage{amsmath, amssymb, amsthm}
\usepackage[margin=1in]{geometry}
\usepackage{hyperref}

\newcommand{\R}{\mathbb{R}}
\newtheorem{theorem}{Theorem}

\title{Notes on mathematical analysis}
\author{A working notebook}
\date{\today}

\begin{document}
\maketitle

\section{Introduction}
These notes explore the connections between continuity,
differentiation, and integration. We begin with a familiar
idea: understanding a function through its local behavior.

\section{Limits and continuity}
\label{sec:limits}
Let $f : \R \to \R$. We say that $f$ is continuous at $a$
if, for every $\varepsilon > 0$, there exists $\delta > 0$ such that
\begin{equation}
  |x-a| < \delta \implies |f(x)-f(a)| < \varepsilon.
  \label{eq:continuity}
\end{equation}

\subsection{A useful example}
Consider the following expression:
\[
  f(x) = \frac{\sin^2(\alpha)}{\sqrt{x^2+1}}
\]
% Try typing fr/ or sq/ inside the math expression below.
\[
  x = 0
\]

\section{The fundamental theorem}
\begin{theorem}
If $f$ is continuous on $[a,b]$, then
\[
  \frac{d}{dx}\int_a^x f(t)\,dt = f(x).
\]
\end{theorem}

\section{Further remarks}
\label{sec:remarks}
The continuity condition in Section~\ref{sec:limits} is essential.
For a more detailed treatment, see \cite{rudin1976}.

\subsection{Next steps}
Explore uniform convergence and the interchange of limits.

\bibliographystyle{plain}
\bibliography{references}
\end{document}
